A finite-change decomposition makes the distinction between insurance and proceeds precise. At a common face $d$, define $A_d=\{d>cg,\ dH_G(g)\le cg\}$ as the states in which the clause prevents default, and $B_d=\{g<\bar g,\ d\le cg\}$ as the states receiving relief even though straight debt would repay. With zero recovery,
\begin{equation}\label{eq:finite_tradeoff}
 b_G(d;c)-b_S(d;c)
 =\frac{d}{R_f}\big[(1-\kappa)\Prob(A_d)-\kappa\Prob(B_d)\big].
\end{equation}
To obtain this identity, partition states according to repayment under the two contracts. States in which both default contribute zero; states outside the trigger in which both survive contribute the same repayment. In $A_d$, the newly surviving claim pays $(1-\kappa)d$. In $B_d$, the clause gives up $\kappa d$. Summing these differences gives \eqref{eq:finite_tradeoff}, without requiring the clause's survival set to be a single interval. The local derivative in \eqref{eq:tradeoff} is the infinitesimal version on the stated threshold branch.

Figure~\ref{fig:loading} illustrates why these terms must be considered jointly. The clause can sharply reduce default at the illustrated face while lowering proceeds, because payments forgone in already solvent states outweigh the discounted value of new repayments. The reduction in default probability alone therefore does not establish an expansion of fiscal space. Increasing relief can enlarge $A_d$, but it also lowers what creditors receive in every newly surviving state and in $B_d$. A useful trigger concentrates the first effect where otherwise collectible payments would be lost to default and limits the second effect where resources were already sufficient.

Reoptimizing issuance changes the comparison. The sovereign can choose a different face under the clause, so the sign of \eqref{eq:finite_tradeoff} at one illustrated face need not equal the sign of the difference between the two maximum proceeds. In the numerical exercise, the GDP clause supports a higher issue near the binding trigger even though it is costly at low common faces. This is the reason for reporting both the rollover-need table and the optimized fiscal-limit table. The common-face comparison describes the payoff trade-off for a given promised service obligation; the optimized comparison asks how far contract choice moves the peak of the proceeds curve.

The same distinction cautions against ranking instruments by their default probabilities at their respective optima. A better repayment mapping can make a larger face attractive, and that larger face can itself change default exposure. A low optimal default rate can be informative about targeting, but neither minimizing default nor maximizing face alone defines MSB. The relevant objective is the market value of the payments that survive, with the pricing rule and resource process held explicit. When recovery or additional premia are introduced, the value attached to each region changes again.
